\documentclass[12pt]{article}
\usepackage{graphicx} % Required for inserting images
\usepackage[margin=0.5in]{geometry}
\usepackage{amsmath, amssymb}
\usepackage{mathrsfs}


\usepackage{soul}


% Dr. David Brown Commands
\newcommand{\ds}{\displaystyle}
\newcommand{\R}{\mathbb{R}}
\newcommand{\N}{\mathbb{N}}
\newcommand{\Q}{\mathbb{Q}}
\newcommand{\C}{\mathbb{C}}


% Commands to get script r
\usepackage{calligra}
\DeclareMathAlphabet{\mathcalligra}{T1}{calligra}{m}{n}
\DeclareFontShape{T1}{calligra}{m}{n}{<->s*[2.2]callig15}{}
\newcommand{\scripty}[1]{\ensuremath{\mathcalligra{#1}}}

% Unit commands
\newcommand{\degree}{\text{°}}
\newcommand{\unitSpeed}{\frac{\text{m}}{\text{s}}}
\newcommand{\unitAcceleration}{\frac{\text{m}}{\text{s}^2}}

% Constant commands
\newcommand{\avogadro}{6.022\times10^{23}}

\newcommand{\boltzmannConstK}{1.381\times10^{-23}\frac{\text{J}}{\text{K}}}
\newcommand{\boltzmannConstInEV}{8.617\times10^{-5}\frac{\text{eV}}{\text{K}}}
\newcommand{\planckConst}{6.626\times10^{-34}\text{J}\cdot\text{s}}

\newcommand{\atmP}{1.01325\times10^5\,\text{Pa}}

% Commands to easily write partial derivatives
\newcommand{\partDeriv}[2]{\frac{\partial #1}{\partial #2}}
\newcommand{\doublePartDeriv}[2]{\frac{\partial^2 #1}{\partial^2 #2}}


\title{Problem 7.22}
\author{Gabriel Decker}


\begin{document}


\maketitle
\begin{flushleft}


In this problem, we will consider a degenerate electron gas where each electron has highly relativistic speed.
This implies that their energies are essentially $\epsilon=|\vec{p}|c$.
Given that, we will derive the chemical potential (or Fermi energy) and the total energy.\\~\\

\textbf{Part a}\\

First, let's begin with the energy $\epsilon$ of the electrons.
It depends on the momentum of the electrons.
We can do a similar thing as the book's derivation, where we assume that our system is a box and only certain de Broglie wavelenghts and momenta are allowed.
$\lambda_n=2L/n$ and $p_n=h/\lambda_n$ imply that $p=hn/2L$.
This yields the following for each direction $q\in[x,y,z]$ for the system.
$$p_q=\frac{hn_q}{2L}$$\\~\\

With this, we can put the three components of momentum into a vector.
$$\vec{p}=\frac{hn_x}{2L}\hat{x}+\frac{hn_y}{2L}\hat{y}+\frac{hn_z}{2L}\hat{z}$$
$$\implies\vec{p}=\frac{h}{2L}\left[n_x\hat{x}+n_y\hat{y}+n_z\hat{z}\right]$$
We can describe a space of $n$s ($n$-space) where each coordinate's corresponding $n$ has its own axis and describe a molecule's state with a vector.
Each point in the grid (with one unit separating points) represents a different state because each point has a different configuration of $n_q$s.
$$\implies\vec{p}=\frac{h}{2L}\vec{n}$$
$$\implies|\vec{p}|=\frac{h}{2L}|\vec{n}|=\frac{h}{2L}\sqrt{n_x^2+n_y^2+n_z^2}$$\\~\\

By plugging in our relation for relativistic energy $\epsilon=|\vec{p}|c$, we get the allowed energy states.
$$\epsilon=|\vec{p}|c$$
$$\implies\epsilon=\frac{hc}{2L}|\vec{n}|$$\\~\\

In our degenerate electron gas, the electrons are in the lowest states.
Therefore, we, like the book, can think of the $n$-space as approximately one eighth of a sphere (since $n_q>0$).
There is also no preference, direction-wise.
So, we can define $n_\text{max}=|\vec{n}|$ for when we're interested in the Fermi energy.
$n_\text{max}$ is the maximum $n$ in a particular direction.
This is because Fermi energy $\epsilon_F=\mu$ is the highest energy level of the system.
This works because the radius of the sphere is approximately the same in all directions, and each direction's max $n$ implies that the the other $n$s are minimal.
Therefore, the fermi temperature (or $\mu$) is the following.
$$\implies\epsilon_F=\frac{hc}{2L}n_\text{max}$$\\~\\

We now want to express $\epsilon_F$ in more general terms.
Specifically, in terms of $N/V$.
$V$ is easy, since it's simply $V=L^3$ in our box.\\~\\

If we consider the eighth of a sphere in $n$-space again, the total volume of the eigth-sphere equals the number of enclosed points.
This works because each state has at least one unit separating itself from every other state, so we can think of each grid point as having a 1-by-1-by-1 cube surrounding it.
This implies that the total number of electrons is two (accounting for spin up and down) times the volume of this eighth-sphere.
This cube has a volume of one, so the enclosed volume thus represents the number of points or states.
Thus, since $n_\text{max}$ is again, the radius of the sphere,
$$N=2\cdot\frac{1}{8}\cdot\frac{4}{3}\pi n_\text{max}^3$$
$$\implies N=\frac{\pi n_\text{max}^3}{3}$$
This is the same as the text's derivation.\\~\\

We can combine these into the $\epsilon_F$ equation by doing the following.
$$\epsilon_F=\frac{hc}{2L}n_\text{max}$$
$$\implies\epsilon_F=\frac{hc}{2}\cdot\left(\frac{3}{\pi}\right)^{1/3}\cdot\left(\frac{\pi}{3}\right)^{1/3}\cdot\frac{n_\text{max}}{L}$$
$$\implies\epsilon_F=\frac{hc}{2}\cdot\left(\frac{3}{\pi}\right)^{1/3}\left[\left[\left(\frac{\pi}{3}\right)^{1/3}\cdot\frac{n_\text{max}}{L}\right]^3\right]^{1/3}$$
$$\implies\epsilon_F=\frac{hc}{2}\cdot\left(\frac{3}{\pi}\right)^{1/3}\left[\frac{\pi n_\text{max}^3}{L^3}\right]^{1/3}$$
$$\implies\epsilon_F=\frac{hc}{2}\cdot\left(\frac{3}{\pi}\right)^{1/3}\left[\frac{N}{V}\right]^{1/3}$$
$$\implies\epsilon_F=hc\left[\frac{3N}{8\pi V}\right]^{1/3}$$\\~\\

\textbf{Part b}\\
Now that we have an equation for the Fermi energy, we can derive an equation for the total energy.
The total energy of our system is simply the sum of the individual energies of our electrons.
Therefore, we get the following.
$$U=2\sum_{n_x}\sum_{n_y}\sum_{n_z}\epsilon(\vec{n})$$
If we assume that we have enough electrons to make a pretty good sphere in $n$-space, we can treat this as an integral.
$$\implies U=2\int\int\int\epsilon(\vec{n})\,dn_xdn_ydn_z$$
We can simplify this by using our $n$-space eight-sphere again and use spherical coordinates.
Since the radius is $n_\text{max}$, we can use that in our integral.
Recall that in spherical coordinates, $dV=n^2\sin(\theta)dnd\theta d\phi$.
Also, in general, $|\vec{n}|=n$.
$$\implies U=2\int\int\int\epsilon(n)n^2\sin(\theta)dnd\theta d\phi$$
$$\implies U=2\int_0^{n_\text{max}}\epsilon(n)n^2\,dn\int_0^{\pi/2}\sin(\theta)\,d\theta \int_0^{\pi/2}\,d\phi$$
$$\implies U=2\int_0^{n_\text{max}}\epsilon(n)n^2\,dn\cdot1\cdot\frac{\pi}{2}$$
$$\implies U=\pi\int_0^{n_\text{max}}\epsilon(n)n^2\,dn$$
We can substitute our expression for $\epsilon$ to continue to solve this.
$$\implies U=\pi\int_0^{n_\text{max}}\frac{hc}{2L}n\cdot n^2\,dn$$
$$\implies U=\frac{hc\pi}{2L}\int_0^{n_\text{max}}n^3\,dn$$
$$\implies U=\frac{hc\pi}{8L}n_\text{max}^4$$\\~\\

Now, expressing $U$ in terms of $N$ and $\epsilon_F$, we get the following.
$$U=\frac{hc\pi}{8L}n_\text{max}^4$$
$$\implies U=\frac{hc}{2L}n_\text{max}\cdot\frac{\pi n_\text{max}^3}{3}\cdot\frac{3}{4}$$
$$\implies U=\frac{3}{4}N\epsilon_F$$



\end{flushleft}

\end{document}
